% context: markscheme (official solutions) for 2-6DN_Mental-math
% topic: mental math, no calculator, answers only — metric length prefixes, squares
%        and square roots, m/s and km/h, distance = speed x time, average speed vs
%        average velocity, fractions to decimals, fractions of an hour, acceleration,
%        free fall with g = 10 m/s^2
% convention: answers in bold; "=" joins exact equalities, "\approx" a rounded or
%        approximate value; italic notes give the one-line reasoning and the usual traps
\documentclass[12pt, twoside]{article}
%\documentclass[12pt, twoside]{article}
\usepackage[letterpaper, margin=1in, headsep=0.2in]{geometry}
\setlength{\headheight}{0.6in}
%\usepackage[english]{babel}
\usepackage[utf8]{inputenc}
\usepackage{microtype}
\usepackage{amsmath}
\usepackage{amssymb}
%\usepackage{amsfonts}
\usepackage[nomessages]{fp} %\FPeval{\var-name}{2*sin(pi/6)}
\usepackage{siunitx} %units in math. eg 20\milli\meter
\usepackage{yhmath} % for arcs, overparenth command
\usepackage{tikz} %graphics
\usetikzlibrary{quotes, angles, arrows, arrows.meta}
\usepackage{graphicx} %consider setting \graphicspath{{images/}}
\usepackage{parskip} %no paragraph indent
\usepackage{enumitem}
\usepackage{multicol}
\usepackage{venndiagram}

\usepackage{fancyhdr}
\pagestyle{fancy}
\fancyhf{}
\renewcommand{\headrulewidth}{0pt} % disable the underline of the header
\raggedbottom
\hfuzz=2mm %suppresses overfull box warnings

\usepackage{hyperref}
\usepackage{float}
\setlength{\parskip}{0.2\baselineskip}

\title{Physics}
\author{Chris Huson}
\date{November 2026}

\fancyhead[LE]{\thepage}
\fancyhead[RO]{\thepage}
\fancyhead[LO]{La Scuola d'Italia / Huson / Physics \\* 5 November 2026}

\begin{document}

\subsubsection*{2.6 Warm-up: Mental Math Solutions}

\emph{Students write only the answer. The notes give the step a student should be
doing in their head, and the mistakes worth looking for.}

\begin{enumerate}[itemsep=0.3cm]

\item Fill in the blanks.
  \begin{multicols}{2}
    \begin{enumerate}[itemsep=0.1cm]
      \item 3.2 km = \textbf{3200 m}
      \item 60 cm = \textbf{0.6 m}
      \item 1.5 m = \textbf{150 cm}
      \item 9 mm = \textbf{0.9 cm}
      \item 0.4 m = \textbf{400 mm}
      \item 750 m = \textbf{0.75 km}
    \end{enumerate}
  \end{multicols}
  \emph{Going to a smaller unit makes the number bigger. For (f), going to a bigger
  unit, divide by 1000: 750 m is three quarters of a kilometer.}

\item Squares and square roots.
  \begin{multicols}{4}
    \begin{enumerate}[itemsep=0.1cm]
      \item $4^2 = \mathbf{16}$
      \item $9^2 = \mathbf{81}$
      \item $11^2 = \mathbf{121}$
      \item $30^2 = \mathbf{900}$
      \item $0.2^2 = \mathbf{0.04}$
      \item $1.5^2 = \mathbf{2.25}$
      \item $\sqrt{49} = \mathbf{7}$
      \item $\sqrt{400} = \mathbf{20}$
    \end{enumerate}
  \end{multicols}
  \emph{The traps: $0.2^2 = 0.2 \times 0.2 = 0.04$, not 0.4, and $30^2 = 900$, not 90.
  For (f), $1.5 \times 1.5 = 1.5 + 0.75$. For (h), $20 \times 20 = 400$.}

\item Speed units.
  \begin{multicols}{2}
    \begin{enumerate}[itemsep=0.1cm]
      \item 5 m/s = \textbf{18 km/h}
      \item 15 m/s = \textbf{54 km/h}
      \item 72 km/h = \textbf{20 m/s}
      \item 90 km/h = \textbf{25 m/s}
    \end{enumerate}
  \end{multicols}
  \emph{From m/s to km/h, multiply by 3.6; back again, divide. 72 is $2 \times 36$,
  and 90 is $2.5 \times 36$. The number in km/h is always the bigger one.}

\item Distance = speed $\times$ time.
  \begin{enumerate}[itemsep=0.1cm]
    \item \textbf{30 m} \quad $1.5 \times 20 = 30$.
    \item \textbf{360 m} \quad One minute is 60 s, and $6 \times 60 = 360$.
    \item \textbf{30 km} \quad 20 minutes is $\frac{1}{3}$ hour, and
      $90 \times \frac{1}{3} = 30$.
    \item \textbf{20 s} \quad time = distance $\div$ speed $= 100 \div 5$.
  \end{enumerate}
  \emph{In (b) and (c) the time has to be in the same unit as the speed first. A
  student who writes 6 m or 1800 km has skipped that step.}

\item One lap of the track.

  Average speed = \textbf{5 m/s} \quad ($400 \div 80$) \qquad
  Average velocity = \textbf{0} \quad (displacement is zero)

  \emph{The runner finishes where the lap began, so the displacement is 0, and so is
  the average velocity, however fast the runner went.}

\newpage

\item Fractions as decimals.
  \begin{multicols}{4}
    \begin{enumerate}[itemsep=0.15cm]
      \item $\dfrac{3}{5} = \mathbf{0.6}$
      \item $\dfrac{9}{10} = \mathbf{0.9}$
      \item $\dfrac{3}{20} = \mathbf{0.15}$
      \item $\dfrac{1}{50} = \mathbf{0.02}$
      \item $\dfrac{3}{25} = \mathbf{0.12}$
      \item $\dfrac{5}{4} = \mathbf{1.25}$
      \item $\dfrac{1}{9} = \mathbf{0.\overline{1}}$
      \item $\dfrac{4}{3} = \mathbf{1.\overline{3}}$
    \end{enumerate}
  \end{multicols}
  \emph{Think in money: $\frac{3}{20}$ is three nickels, $\frac{1}{50}$ is 2 cents,
  $\frac{3}{25}$ is 12 cents. Accept $0.111\ldots$ and $1.333\ldots$, or 0.11 and
  1.33 written with $\approx$.}

\item Fractions of an hour and a week.
  \begin{multicols}{2}
    \begin{enumerate}[itemsep=0.1cm]
      \item 0.5 hour = \textbf{30 min}
      \item 0.4 hour = \textbf{24 min}
      \item 2.5 min = \textbf{150 s}
      \item 90 s = \textbf{1.5 min}
      \item 6 min = \textbf{0.1 h}
      \item 40 min = $\mathbf{\dfrac{2}{3}}$ \textbf{h}
      \item 1 week = \textbf{168 hours}
      \item 2.25 hours = \textbf{2 h 15 min}
    \end{enumerate}
  \end{multicols}
  \emph{Each tenth of an hour is 6 minutes. For (g), $7 \times 24 = 140 + 28$. In (h),
  0.25 h is 15 min, not 25 min.}

\item Acceleration = change in velocity $\div$ time.
  \begin{enumerate}[itemsep=0.1cm]
    \item \textbf{4 m/s$^\mathbf{2}$} \quad $20 \div 5$.
    \item \textbf{2 m/s$^\mathbf{2}$} \quad The change is $8 - 2 = 6$ m/s, and
      $6 \div 3 = 2$. \emph{The slip is $8 \div 3$, using the final speed instead of
      the change.}
    \item \textbf{4 s} \quad $20 \div 5$.
    \item Average velocity = \textbf{10 m/s} \quad ($\frac{0 + 20}{2}$); \quad
      distance = \textbf{50 m} \quad ($10 \times 5$).
  \end{enumerate}

\item Falling with $g \approx 10$ m/s$^2$.
  \begin{enumerate}[itemsep=0.1cm]
    \item \textbf{30 m/s} \quad $10 \times 3$.
    \item \textbf{20 m} \quad $\frac{1}{2} \times 10 \times 2^2 = 5 \times 4$.
      \emph{Not 10 m: the stone is only going 20 m/s at the end, so its average is
      10 m/s.}
    \item \textbf{2 s} \quad It loses 10 m/s every second, and $20 \div 10 = 2$.
    \item Velocity = \textbf{0}; \quad acceleration = \textbf{10 m/s$^\mathbf{2}$
      downward}.

      \emph{The common wrong answer is $a = 0$ at the top. Gravity does not switch
      off; the ball is still slowing down, then speeding up downward.}
  \end{enumerate}

\item \textbf{About 26 miles.} \quad $42 \times \dfrac{5}{8} = \dfrac{210}{8}
  = 26\dfrac{1}{4}$.

  \emph{Accept 25 to 26.25 mi. The marathon is 42.195 km, which is 26.2 mi.}

\end{enumerate}
\end{document}
